\documentclass[11pt,a4paper]{article}
\usepackage[margin=20mm]{geometry}
\usepackage{fontspec}
\setmainfont{DejaVu Serif}
\setsansfont{DejaVu Sans}
\setmonofont{DejaVu Sans Mono}[Scale=0.85]
\usepackage{polyglossia}
\setdefaultlanguage{russian}
\usepackage{graphicx}
\usepackage{hyperref}
\usepackage{enumitem}
\usepackage{fancyvrb}
\hypersetup{hidelinks}
\setlength{\parindent}{0pt}
\setlength{\parskip}{4pt}
\usepackage{titlesec}
\titleformat{\section}{\large\bfseries}{}{0pt}{}
\titlespacing*{\section}{0pt}{10pt}{4pt}
\sloppy
\setlist{nosep,leftmargin=*}
\emergencystretch=3em
\newcommand{\shot}[2]{\begin{center}\includegraphics[width=\linewidth,height=.34\textheight,keepaspectratio]{#1}\par\small #2\end{center}}
\begin{document}
\begin{center}
{\small Университет ИТМО · Фронт-энд разработка · 2026}\par
{\LARGE Лабораторная работа № 2}\par
{\large ML-платформа с JSON API}\par
Светлов Ярослав, группа J3211
\end{center}
\section*{Цель и технологии}
Подключить интерфейс ML-платформы к серверу данных. Использованы HTML, Bootstrap, обычный JavaScript, Chart.js, JSON Server 0.17.4 и json-server-auth 2.1.0. Vue не используется.
\section*{Реализованные возможности}
\begin{itemize}
\item Регистрация, вход, восстановление сессии после обновления страницы и выход. Недействительный токен возвращает пользователя на вход.
\item Кабинет показывает профиль, собственные эксперименты, модели и артефакты.
\item Серверный поиск по дате, минимальной accuracy, максимальной latency и тексту; сброс и пустая выдача.
\item Детали по ID: метрики, график, логи и скачивание JSON/TXT.
\item Добавление версии со связью со своим экспериментом; статусы Staging, Production, Archived сохраняются сервером.
\end{itemize}

\section*{Организация кода и доступ}
Статика находится в \texttt{public}, сервер запускается через \texttt{server.js}. Один помощник fetch добавляет Bearer-токен и обрабатывает ответы. Интерфейс сообщает о загрузке, ошибках и пустых данных.

Используются маршруты \texttt{/register}, \texttt{/login}, \texttt{/users/:id}, \texttt{/experiments}, \texttt{/experiments/:id}, \texttt{/models}. Сервер разрешает только свой профиль и изменение собственных моделей. Фильтры передаются как \texttt{date}, \texttt{accuracy\_gte}, \texttt{latencyMs\_lte}, \texttt{q}.
\section*{Результаты проверки}
19.09.2026 выполнен повторный проход в Firefox 148.0.2 и Chromium 146.0.7680.153: 26 сценариев Firefox и 6 дополнительных сценариев Chromium пройдены. Проверены регистрация, повторный email, неверный пароль, восстановление и завершение сессии, отдельные и совместные фильтры, сброс и пустой результат, четыре разных эксперимента, неизвестный ID, скачивание JSON/TXT, создание версии и смена статусов с перезагрузкой. Проверены недействительный токен, недоступный API и узкие окна (Firefox 500 px, Chromium 390 px). Девять API-тестов пройдены; среди них проверки отсутствия токена, чужой модели, содержимого файлов и сохранения на диск.
\section*{Запуск и проверка}
Требуется Node.js 22.12 или новее. Команды выполняются из папки этой работы.
\begin{verbatim}
npm ci
npm start
\end{verbatim}
Открыть \url{http://localhost:3000}. Автотесты: \texttt{npm test}.
Учебные аккаунты: \texttt{admin@nexus.ai / password} и \texttt{vision@nexus.ai / cv12345}. Данные сохраняются в \texttt{db.local.json}, исходная база \texttt{db.json} не меняется.\section*{Вывод и ограничения}
Обучение и облачное развёртывание не выполняются: статус Production имитирует деплой. Артефакты содержат демонстрационные метрики и логи, а не веса моделей. JSON Server с json-server-auth предназначен только для localhost и учебных данных; нельзя использовать реальные пароли.

\section*{Скриншоты}
\shot{checks/lab2-experiment.png}{Рисунок 1. Метрики и график выбранного эксперимента.}
\shot{checks/lab2-combined-filters.png}{Рисунок 2. Совместная серверная фильтрация.}

\end{document}
